\SourcePage{311}
\SGAsection{5}{共轭定理}
\begin{SGAresult}{定理}{5.1}
设 $S$ 是\Term{semilocal-scheme}{半局部概形}，$G$ 是\Term{reductive-group}{约化 $S$ 群}，$P$ 和 $P'$ 是 $G$ 的两个\Term{opposite-parabolic}{相反抛物子群}（4.3.3）。则
\[
\operatorname{rad}^{u}(P)(S)\,P'\cdot P=G,
\]
即当 $u$ 遍历 $\operatorname{rad}^{u}(P)(S)$ 时，各\Term{open-subset}{开集} $uP'\cdot P$（4.3.6）的并是整个 $G$。
\end{SGAresult}
证明分成以下几步。

\textbf{5.1.1。}只需在 $S$ 是一个域 $k$ 的谱的情形下证明；这立即由 2.6 得到。

\textbf{5.1.2。}设 $L=P\cap P'$。假设 $L$ 有一个\Term{borel-subgroup}{Borel 子群} $B_L$；取 $B_L$ 的一个\Term{maximal-torus}{极大环面} $T$（第 XXII 讲 5.9.7）。立即可验 $B=B_L\cdot\operatorname{rad}^{u}(P)$ 是 $P$ 的一个\Term{borel-subgroup}{Borel 子群}；设 $B'$ 是 $G$ 中相对于 $T$ 与 $B$ 相反的\Term{borel-subgroup}{Borel 子群}（即 $B\cap B'=T$）。将 $G$ 相对于 $T$ \Term{split-group}{分裂}后立即可验 $B'\subset P'$。我们证明
\[
 B^u(S)\,B'\cdot B\subset\operatorname{rad}^{u}(P)(S)\,P'\cdot P.\tag{x}
\]
因为
\[
 B^u(S)\subset P(S)=\operatorname{rad}^{u}(P)(S)\cdot L(S)
 \subset\operatorname{rad}^{u}(P)(S)P'(S),
\]
只需证明 $B'\cdot B\subset P'\cdot P$，这是显然的。由（x）可知，只需对 $(B,B')$ 证明 5.1。

\textbf{5.1.3。}若 $k$ 是\Term{algebraically-closed-field}{代数闭域}，则定理成立：事实上，5.1.2 的条件得到满足，由第 XXII 讲 5.7.10 即可得出结论。

\textbf{5.1.4。}当 $k$ 是无限域时，定理成立。事实上，由 2.7，$\operatorname{rad}^{u}(P)(k)$ 在 $\operatorname{rad}^{u}(P)(\bar k)$ 中\Term{dense}{稠密}，而定理对 $\bar k$ 成立。

\textbf{5.1.5。}于是归结为 $k$ 是\Term{finite-field}{有限域}的情形。$\operatorname{Bor}(L)$ 是 $L$ 的\Term{smooth}{光滑}\Term{homogeneous-space}{齐性空间}；因此由 Lang 定理（\textit{Am. J. of Maths.}, 78, 1956\footnote{编者注：也见 [DG70]，\S III.5，7.4。}），$L$ 有一个\Term{borel-subgroup}{Borel 子群} $B_L$。由 5.1.2，可以假设 $P=B$ 和 $P'=B'$ 都是\Term{borel-subgroup}{Borel 子群}。记 $T=B\cap B'$。

\SourcePage{312}
\textbf{5.1.6。}设 $K$ 是 $k$ 的\Term{algebraic-closure}{代数闭包}；为三元组 $(G_K,B_K,T_K)$ 选取一个\Term{pinning}{钉扎}，以 $R_+$（分别以 $\Delta$）表示\Term{positive-root}{正根}（分别为\Term{simple-root}{单根}）的\Term{set}{集合}。由第 XXII 讲 5.7.2，只需证明：对任意 $\alpha\in\Delta$，有
\[
 u_\alpha B'^u(K)\subset B^u(k)B'^u(K)B(K).\tag{1}
\]
设 $\alpha_i$ 是 $\alpha$ 在 $k$ 上的各个\Term{conjugacy}{共轭}\Term{root}{根}（因 $B$“定义在 $k$ 上”，它们都属于 $\Delta$），设 $R'$ 是可写成 $\alpha_i$ 的线性组合的\Term{root}{根}的\Term{set}{集合}。记 $R'_- =R_-\cap R'$。因“$R'$ 定义在 $k$ 上”，$T$ 有一个\Term{subtorus}{子环面} $Q$，使 $Q_K$ 是各 $\alpha_i$ 的公共核的\Term{maximal-torus}{极大环面}。

记 $Z=\operatorname{Centr}_G(Q)$、$B_Z=B\cap Z$、$B'_Z=B'\cap Z$（见第 XXII 讲 5.10.2）。我们说明只需在 $Z$ 中验证所求断言，即
\[
 u_\alpha\cdot B'^u_Z(K)\subset B^u_Z(k)B'^u_Z(K)B_Z(K).\tag{2}
\]
有 $(B'^u_Z)_K=\prod_{\alpha\in R'_-}U_\alpha$。设 $R''$ 是 $R'$ 在 $R$ 中的补集，置
\[
 V=\prod_{\alpha\in R''\cap R_-}U_\alpha.
\]
立即有 $(B'^u)_K=(B'^u_Z)_K\cdot V$，且 $(B_Z)_K$ \Term{normalization}{正规化} $V$（第 XXII 讲 5.6.7）。因此由（2）依次得到
\begin{align*}
 u_\alpha\cdot B'^u(K)
 &=u_\alpha\cdot B'^u_Z(K)V(K)\\
 &\subset B^u_Z(k)B'^u_Z(K)B_Z(K)V(K)\\
 &\subset B^u_Z(k)B'^u_Z(K)V(K)B_Z(K),
\end{align*}
立即推出（1）。于是归结为 $G=Z$ 的情形，即 $K/k$ 的 Galois 群在\Term{simple-root}{单根}上\Term{transitive-action}{传递作用}的情形。

\textbf{5.1.7。}待证断言\Term{equivalence}{等价}于：$G/B$ 是 $B'^u$ 的像所成\Term{open-subset}{开集}在 $B^u(k)$ 下的各\Term{translation}{平移}的并。将 $G$ 换成它的\Term{adjoint-group}{伴随群}（甚至换成任何具有同一\Term{adjoint-group}{伴随群}的群），这个断言不变。因此可以假设 $G$ 是\Term{adjoint-group}{伴随群}。

\textbf{5.1.8。}现在考虑 $G_K$ 的 Dynkin 图。Galois 群在此 Dynkin 图上\Term{transitive-action}{传递作用}。但这个 Galois 群只有循环\Term{quotient-group}{商群}，而 Dynkin 图没有圈。因此立即可知，此图为 $nA_1$ 型，$n\geq0$，或 $mA_2$ 型，$m\geq0$。利用第 XXIV 讲 5.9 的\Term{canonical}{典范}\Term{resolution}{分解}，可写成
\[
 G=\prod_{D/K}G_0,
\]
其中 $D$ 是\Term{finite}{有限} $K$ \Term{scheme}{概形}，而 $G_0$ 或是\Term{torus}{环面}，或为 $A_1$ 型，或为 $A_2$ 型。
由第 XXIV 讲 5.12，$B$ 来自 $G_0$ 的一个\Term{borel-subgroup}{Borel 子群} $B_0$，$T$ 来自 $G_0$ 的一个\Term{maximal-torus}{极大环面} $T_0$；$B'$ 来自 $G_0$ 中相对于 $T_0$ 与 $B_0$ 相反的\Term{borel-subgroup}{Borel 子群} $B'_0$。
\SourcePage{313}
有
\[
 B'^u(k)=B_0^u(D),\qquad
 B'^u\cdot T\cdot B^u=\prod_{D/k}B'^u_0\cdot T_0\cdot B^u_0,
\]
所以只需对三元组 $(G_0,B_0,T_0)$ 证明所求断言。

\textbf{5.1.9。}因此可以假设 $G$ 为 $\varnothing$、$A_1$ 或 $A_2$ 型。因为 $G$ 有一个\Term{borel-subgroup}{Borel 子群} $B$，$G$ 相对于 $B$ 是\Term{quasi-split-group}{拟分裂的}（第 XXIV 讲 3.9.1），所以若它为 $\varnothing$ 型或 $A_1$ 型，就可\Term{split-group}{分裂}。定理已在\Term{split-group}{分裂}情形得到证明（第 XXII 讲 5.7.10），因此只需处理 $A_2$ 型。由第 XXIV 讲 3.11，存在\Term{morphism}{态射} $E\to\operatorname{Spec}(k)$，它是 $A_2$ 型 Dynkin 图的\Term{automorphism-group}{自同构群} $\mathbb Z/2\mathbb Z$ 下的 Galois \Term{torsor}{主丛}，使
$G=G^{q\mathrm{\text{É}p}}_{E/\operatorname{Spec}(k)}(A_2)$。
若 $E$ 有\Term{section}{截面}，则 $G$ 可\Term{split-group}{分裂}，定理已经得证。否则必有 $E=\operatorname{Spec}(k')$，其中 $k'$ 是 $k$ 的二次\Term{group-extension}{扩张}。最后，由 5.1.7 中所见，可以假设 $G$ \Term{simply-connected}{单连通}（即 $G$ 是 $\mathrm{SL}_{3,k}$ 的一个\Term{form}{形式}）。

\textbf{5.1.10。}于是处于如下情形：有一个\Term{finite-field}{有限域} $k$ 及其二次\Term{group-extension}{扩张} $k'$。行列式为 $1$ 的 $3\times3$ 矩阵群 $\mathrm{SL}_{3,k'}$ 配备如下\Term{pinning}{钉扎}：\Term{maximal-torus}{极大环面}是对角矩阵群，\Term{borel-subgroup}{Borel 子群}是上三角矩阵群，“钉”是元素
\[
 u_\alpha=\begin{pmatrix}1&1&0\\0&1&0\\0&0&1\end{pmatrix},\qquad
 u_\beta=\begin{pmatrix}1&0&0\\0&1&1\\0&0&1\end{pmatrix}.
\]
立即可见，\Term{big-cell}{大胞腔} $\Omega$ 由下式定义：
\[
 \begin{pmatrix}a&b&c\\d&e&f\\g&h&i\end{pmatrix}\in\Omega(S)
 \quad\Longleftrightarrow\quad a\text{ 和 }ae-bd\text{ 可逆},
\]
并且
\[
 B^u(k)=\left\{\left.\begin{pmatrix}1&\bar x&z\\0&1&x\\0&0&1\end{pmatrix}\ \right|\ x,z\in k',\ z+\bar z=x\bar x\right\}.
\]
要证明 5.1.6 的包含式（1），即要说明：对任意 $a,b,c\in K$（$k$ 的\Term{algebraic-closure}{代数闭包}），存在 $x,z\in k'$，使 $z+\bar z=x\bar x$，且
\[
 \begin{pmatrix}1&\bar x&z\\0&1&x\\0&0&1\end{pmatrix}
 \begin{pmatrix}1&1&0\\0&1&0\\0&0&1\end{pmatrix}
 \begin{pmatrix}1&0&0\\a&1&0\\b&c&1\end{pmatrix}
 \subset\Omega(K).\tag{1}
\]
若 $a\neq-1$，取 $x=z=0$。
\SourcePage{314}
若 $a=-1$，须满足的条件写为
\[
 \begin{cases}
 z+\bar z=x\bar x,\\
 bz-\bar x\neq0,\\
 (b+c)\bar z+bx-1\neq0.
 \end{cases}
\]
设 $q$ 是 $k$ 的元素个数（$q\geq2$）。已知对任意 $m\in k$，方程 $z+\bar z=m$ 在 $k'$ 中有 $q$ 个解。
若 $b=0$，取 $x=1$；须解 $z+\bar z=1$、$c\bar z\neq1$，由上面的说明总能做到。
若 $b\neq0$，取 $x=0$；须解
\[
 z+\bar z=0,\qquad z\neq0,\qquad(b+c)\bar z\neq1.
\]
若 $q\geq3$，总能做到。若 $k=\mathbb F_2$，当 $b+c\neq1$ 时也能做到，可以取 $z=1$。
所以只剩下 $k=\mathbb F_2$、$b+c=1$、$b\neq0$ 的情形。方程组这时写为
\[
 z+\bar z=x\bar x,\qquad bz\neq\bar x,\qquad\bar z+bx\neq1.
\]
若 $b=1$（或 $b\notin k'$），取 $x=1$；后两个条件写为 $z\neq b^{-1},1-b$，它们都是 $z+\bar z=1$ 的推论，而该方程有解。最后，若 $b\in k'-k$，可以取 $x=\bar b$、$z=b$。证毕。

\begin{SGAresult}{推论}{5.2}
设 $S$ 是\Term{semilocal-scheme}{半局部概形}，$G$ 是\Term{reductive-group}{约化 $S$ 群}，$P$ 和 $P'$ 是 $G$ 的两个\Term{opposite-parabolic}{相反抛物子群}。\Term{canonical}{典范}映射
\[
 \operatorname{rad}^{u}(P)(S)\cdot\operatorname{rad}^{u}(P')(S)\longrightarrow(G/P)(S)
\]
是\Term{surjectivity}{满射}（特别地，$(G/P)(S)=G(S)/P(S)$）。$G$ 中与 $P$ 同型的每个\Term{parabolic-subgroup}{抛物子群} $Q$ 都可写成 $\operatorname{int}(uu')P$，其中 $u\in\operatorname{rad}^{u}(P)(S)$、$u'\in\operatorname{rad}^{u}(P')(S)$。
\end{SGAresult}
第二个断言显然\Term{equivalence}{等价}于第一个，我们证明后者。设 $s_i$ 是 $S$ 的\Term{closed-point}{闭点}，$V$ 是 $G/P$ 中 $P'$ 的像所成的\Term{open-subset}{开集}（它\Term{isomorphism}{同构}于 $\operatorname{rad}^{u}(P')$，见 4.3.6），且 $x\in(G/P)(S)$。由 5.1，对每个 $i$，存在\Term{section}{截面} $u_i\in\operatorname{rad}^{u}(P)(\kappa(s_i))$，使 $u_ix_{s_i}$ 是 $V_{s_i}$ 的\Term{section}{截面}。若 $u\in\operatorname{rad}^{u}(P)(S)$ \Term{lifting}{提升}各 $u_i$（2.6），则 $ux$ 是 $V$ 的\Term{section}{截面}，因为此断言可在\Term{closed-fiber}{闭纤维}上验证。但 $\operatorname{rad}^{u}(P')(S)\xrightarrow{\sim}V(S)$ 是\Term{bijection}{双射}，立即得出结论。

\begin{SGAresult}{推论}{5.3}
设 $S$ 是\Term{semilocal-scheme}{半局部概形}，$G$ 是\Term{reductive-group}{约化 $S$ 群}，$P$、$P'$ 和 $Q$ 是 $G$ 的三个\Term{parabolic-subgroup}{抛物子群}。存在 $g\in G(S)$，使 $\operatorname{int}(g)Q$ 相对于 $P$ 和 $P'$ 都处于\Term{transversal-position}{横截位置}。
\end{SGAresult}
用 4.2.4（ii）的记号，须验证 $G$ 中在 $S$ 上\Term{universally-schematically-dense}{普遍概形稠密}的\footnote{编者注：已将“相对\Term{dense}{稠密}”替换为“在 $S$ 上\Term{universally-schematically-dense}{普遍概形稠密}”，见 EGA IV$_3$，定义 11.10.8。}\Term{open-subset}{开集} $\operatorname{Gen}(Q/P)\cap\operatorname{Gen}(Q/P')$ 在 $S$ 上有\Term{section}{截面}。事实上，选取 $G$ 中一个与 $Q$ 相反的\Term{parabolic-subgroup}{抛物子群} $Q_1$（4.3.5（i）），
\SourcePage{315}
置 $U=\operatorname{rad}^{u}(Q)$、$U'=\operatorname{rad}^{u}(Q_1)$。我们证明存在满足要求的 $g\in U(S)U'(S)$。以此\Term{form}{形式}，由 4.2.4（i）和 2.6，只需在 $S$ 各\Term{closed-point}{闭点}上的\Term{fiber}{纤维}中验证，故可假设 $S$ 是一个域 $k$ 的谱。
若 $k$ \Term{algebraically-closed}{代数闭}，存在满足要求的 $g\in G(k)$；由 5.2，$g=uu'q$，其中 $u\in U(k)$、$u'\in U'(k)$、$q\in Q(k)$，且 $\operatorname{int}(uu')Q=\operatorname{int}(g)Q$。
若 $k$ 无限，考虑由如下\Term{cartesian-diagram}{笛卡尔图}定义的 $U\times_kU'$ 的\Term{open-subset}{开集} $V$：
\[
\begin{tikzcd}
 U\times_kU'\arrow[r]&G\\
 V\arrow[u,hook]\arrow[r]&\operatorname{Gen}(Q/P)\cap\operatorname{Gen}(Q/P')\arrow[u,hook].
\end{tikzcd}
\]
由刚才所见，$V(\bar k)\neq\varnothing$，因此 $V$ 在 $U\times_kU'$ 中\Term{dense}{稠密}，由 2.7 它有\Term{section}{截面}。
若 $k$ \Term{finite}{有限}，由 Lang 定理（见 5.1.5），$P$（分别为 $P'$）有一个\Term{borel-subgroup}{Borel 子群} $B$（分别为 $B'$），因为\Term{scheme}{概形} $\operatorname{Bor}(P)\simeq\operatorname{Bor}(P/\operatorname{rad}^{u}(P))$ 和 $\operatorname{Bor}(Q)\simeq\operatorname{Bor}(Q/\operatorname{rad}^{u}(Q))$ 是\Term{smooth}{光滑的}。若 $B_1$ 是与 $B$ 相反的\Term{borel-subgroup}{Borel 子群}（4.3.5（i）），由 5.2 存在 $a\in B^u(k)$、$a_1\in B_1^u(k)$，使 $\operatorname{int}(aa_1)B=B'$；于是 $B_0=\operatorname{int}(aa_1)B_1=\operatorname{int}(a)B_1$ 与 $B'$ 和 $B$ 都相反。若 $Q_0$ 是 $G$ 中唯一包含 $B_0$ 且与 $Q$ 同型的\Term{parabolic-subgroup}{抛物子群}（3.8），则 $Q_0$ 相对于 $P$ 和 $P'$ 处于\Term{transversal-position}{横截位置}（4.2.1（vi））。另一方面，由 5.2，$Q_0$ 可写成 $\operatorname{int}(uu')Q$，其中 $u'\in U'(k)$、$u\in U(k)$，这正是所求。

\begin{SGAresult}{推论}{5.4}
设 $S$ 是\Term{semilocal-scheme}{半局部概形}，$G$ 是\Term{reductive-group}{约化 $S$ 群}，$P$ 和 $Q$ 是 $G$ 的两个\Term{parabolic-subgroup}{抛物子群}。存在 $g\in G(S)$，使 $\operatorname{int}(g)P$ 相对于 $Q$ 处于\Term{osculating-position}{密切位置}，即（4.4.2）$\operatorname{int}(g)P\cap Q$ 是 $G$ 的\Term{parabolic-subgroup}{抛物子群}。
\end{SGAresult}
事实上，由 4.3.5（i），存在 $G$ 中与 $P$ 相反的\Term{parabolic-subgroup}{抛物子群} $P'$。由 5.3，存在 $G$ 中与 $P'$ 同型、相对于 $P$ 和 $Q$ 都处于\Term{transversal-position}{横截位置}的\Term{parabolic-subgroup}{抛物子群} $P'_1$。取 $P'_1\cap Q$ 的一个\Term{maximal-torus}{极大环面} $T$（4.2.7（ii）），设 $P_1$ 是相对于 $T$ 与 $P'_1$ 相反的\Term{subgroup}{子群}，则由 4.4.5，$P_1$ 和 $Q$ 处于\Term{osculating-position}{密切位置}。另一方面，$P$ 和 $P_1$ 都与 $P'_1$ 相反，故存在 $g\in\operatorname{rad}^{u}(P'_1)(S)$，使 $\operatorname{int}(g)P=P_1$（4.3.5（i））。证毕。
还应指出，由同样理由，存在 $u\in\operatorname{rad}^{u}(P)(S)$，使 $\operatorname{int}(u)P'=P'_1$，故 $g$ 可写成 $\operatorname{int}(u)u'$，其中 $u'\in\operatorname{rad}^{u}(P')(S)$；因而
$P_1=\operatorname{int}(uu'u^{-1})P=\operatorname{int}(uu')P$，顺便再次证明了 5.2。
断言 5.3 和 5.4 是本节的主要结果。先列出 5.4 的几个推论。

\begin{SGAresult}{推论}{5.5}
设 $S$ 是\Term{semilocal-scheme}{半局部概形}，$G$ 是\Term{reductive-group}{约化 $S$ 群}。
（i）若 $P$ 和 $Q$ 是 $G$ 的两个\Term{parabolic-subgroup}{抛物子群}，且 $t(P)\subset t(Q)$（见 3.3），则存在 $g\in G(S)$，使 $\operatorname{int}(g)P\subset Q$。
\SourcePage{316}
（ii）设
\[
 P_1\supset P_2\supset\cdots\supset P_n,\qquad
 P'_1\supset P'_2\supset\cdots\supset P'_n
\]
是 $G$ 的两条\Term{parabolic-subgroup}{抛物子群}链，且 $t(P_i)=t(P'_i)$。存在 $g\in G(S)$，使对每个 $i$ 都有 $\operatorname{int}(g)P_i=P'_i$。

（iii）设 $P,Q,P',Q'$ 是 $G$ 的四个\Term{parabolic-subgroup}{抛物子群}，且 $t(P)=t(P')$、$t(Q)=t(Q')$。若对子 $(P,P')$ 和 $(Q,Q')$ 处于\Term{transversal-position}{横截位置}（分别为\Term{osculating-position}{密切位置}），则存在 $g\in G(S)$，使 $\operatorname{int}(g)P=P'$、$\operatorname{int}(g)Q=Q'$。

（iv）设 $P$ 和 $P'$ 是两个同型的\Term{parabolic-subgroup}{抛物子群}，$L$（分别为 $L'$）是 $P$（分别为 $P'$）的一个\Term{levi-subgroup}{Levi 子群}。存在 $g\in G(S)$，使 $\operatorname{int}(g)P=P'$、$\operatorname{int}(g)L=L'$。
\end{SGAresult}
\begin{proof}
（i）立即由 5.4 得到。（ii）对 $n$ 归纳，$n=0$ 时显然；故可假设 $i=1,\ldots,n-1$ 时 $P_i=P'_i$。由 5.2，存在 $g\in G(S)$，使 $\operatorname{int}(g)P_n=P'_n$；但这时 $P_n$ 和 $\operatorname{int}(g)P_n$ 都包含在 $P_{n-1}=P'_{n-1}$ 中，故 $g\in P_{n-1}(S)$（4.4.4（ii）），并且 $i=1,\ldots,n$ 时 $\operatorname{int}(g)P_i=P'_i$。
另一方面，（iv）立即由 5.2 和 1.8 得到。证明（iii）的“\Term{transversal-position}{横截位置}”情形；当 $P$ 和 $Q$ 的型相反（4.3.3（iii））时，断言是（iv）的推论。一般情形下，由 4.2.7（iii）和 4.4.6，可以分别在 $P,P',Q,Q'$ 中找到\Term{parabolic-subgroup}{抛物子群} $P_1,P'_1,Q_1,Q'_1$，使 $P_1$ 和 $P'_1$ 相反，$Q_1$ 和 $Q'_1$ 也相反，且 $t(P_1)=t(P'_1)$；因此存在 $g\in G(S)$，使 $\operatorname{int}(g)P_1=P'_1$、$\operatorname{int}(g)Q_1=Q'_1$，于是可假设 $P_1=P'_1$、$Q_1=Q'_1$。但这时 $P$ 和 $P'$ 处于\Term{osculating-position}{密切位置}且同型，故 $P=P'$（4.4.4（i））；同理 $Q=Q'$。
最后证明（iii）的“\Term{osculating-position}{密切位置}”情形。由\Term{conjugacy}{共轭}定理（5.2），可假设 $P=P'$；仍由该定理，可找到 $g\in G(S)$，使 $\operatorname{int}(g)Q=Q'$；但由 4.4.4（ii），$g\in P(S)$，故 $\operatorname{int}(g)P=P=P'$。
\end{proof}

\begin{SGAresult}{定义}{5.6}
设 $S$ 是\Term{scheme}{概形}，$G$ 是\Term{reductive-group}{约化 $S$ 群}，$P$ 是 $G$ 的\Term{parabolic-subgroup}{抛物子群}。若 $G$ 的\Term{parabolic-subgroup}{抛物子群} $Q$ 每当包含在 $P$ 中时必有 $Q=P$，就称 $P$ 是\Term{minimal-parabolic}{极小的}。
\end{SGAresult}
须注意，这个概念一般不在取\Term{fiber}{纤维}时保持不变。
\begin{SGAresult}{推论}{5.7}
设 $S$ 是\Term{semilocal-scheme}{半局部概形}，$G$ 是\Term{reductive-group}{约化 $S$ 群}。
（i）设 $t,t'\in\operatorname{Of}(\operatorname{Dyn}(G))(S)$。若 $G$ 中存在 $t$ 型和 $t'$ 型的\Term{parabolic-subgroup}{抛物子群}，则存在 $t\cap t'$ 型的\Term{parabolic-subgroup}{抛物子群}。特别地，当 $P$ 遍历 $G$ 的所有\Term{parabolic-subgroup}{抛物子群}时，各 $t(P)$ 所成\Term{set}{集合}中有一个最小元素 $t_{\min}$。
（ii）$G$ 的每个\Term{parabolic-subgroup}{抛物子群}都包含一个\Term{minimal-parabolic}{极小抛物子群}。$G$ 的\Term{parabolic-subgroup}{抛物子群}是\Term{minimal-parabolic}{极小的}，当且仅当它为 $t_{\min}$ 型。$G$ 的两个\Term{minimal-parabolic}{极小抛物子群}由 $G(S)$ 的一个元素\Term{conjugacy}{共轭}。
\end{SGAresult}
这立即由 5.4 和 5.5（i）得到。

\SourcePage{317}
\textbf{注记 5.8。}与一个\Term{minimal-parabolic}{极小抛物子群}相反的\Term{parabolic-subgroup}{抛物子群}也是\Term{minimal-parabolic}{极小的}；由此推出 $s(t_{\min})=t_{\min}$。
\begin{SGAresult}{推论}{5.9}
设 $S$ 是\Term{scheme}{概形}，$G$ 是\Term{reductive-group}{约化 $S$ 群}，$P$ 是 $G$ 的\Term{parabolic-subgroup}{抛物子群}。\Term{canonical}{典范}\Term{morphism}{态射}
\[
 G\longrightarrow G/P=X
\]
使 $G$ 成为群为 $P_X$ 的、在 Zariski 意义下\Term{locally-trivial-bundle}{局部平凡的 $X$ 丛}。若 $L$ 是 $P$ 的一个\Term{levi-subgroup}{Levi 子群}，则\Term{canonical}{典范}\Term{morphism}{态射}（见 3.12）
\[
 G\longrightarrow G/L=Y
\]
使 $G$ 成为群为 $L_Y$ 的、在 Zariski 意义下\Term{locally-trivial-bundle}{局部平凡的 $Y$ 丛}。
\end{SGAresult}
只需证明：若有\Term{morphism}{态射} $S'\to S$，其中 $S'$ 是\Term{local-scheme}{局部}的，以及\Term{morphism}{态射} $S'\to X$（分别为 $S'\to Y$），就能将后者\Term{lifting}{提升}成\Term{morphism}{态射} $S'\to G$。换言之，可以假设 $S$ \Term{local-scheme}{局部}，而须证明映射 $G(S)\to X(S)$（分别为 $G(S)\to Y(S)$）\Term{surjectivity}{满射}。
第一个断言已在 5.2 中证明；证明第二个。设 $y\in Y(S)$，它在 $X(S)$ 中的\Term{canonical}{典范}像来自某个 $g\in G(S)$；故 $g$ 在 $Y(S)$ 中的\Term{projection}{投影} $y'$ 与 $y$ 在 $X(S)$ 中有相同\Term{projection}{投影}。因此存在唯一的 $u\in\operatorname{rad}^{u}(P)(S)$，使 $y'u=y$，而 $gu$ 在 $Y(S)$ 中的\Term{projection}{投影}确为 $y$。
\begin{SGAresult}{推论}{5.10}
设 $S$ 是\Term{semilocal-scheme}{半局部概形}，$G$ 是\Term{reductive-group}{约化 $S$ 群}。
（i）设 $P$ 是 $G$ 的\Term{parabolic-subgroup}{抛物子群}，$L$ 是 $P$ 的\Term{levi-subgroup}{Levi 子群}。\Term{canonical}{典范}映射（见 3.21）诱导\Term{bijection}{双射}
\[
 H^1(S,L)\xrightarrow{\sim}H^1(S,P)\xrightarrow{\sim}H^1_{t(P)}(S,G).
\]
（ii）设 $t,t'\in\operatorname{Of}(\operatorname{Dyn}(G))(S)$，则（见 3.21）
\[
 H^1_t(S,G)\cap H^1_{t'}(S,G)=H^1_{t\cap t'}(S,G).
\]
（iii）若 $P$ 和 $Q$ 是 $G$ 的两个处于\Term{osculating-position}{密切位置}的\Term{parabolic-subgroup}{抛物子群}，则以下\Term{canonical}{典范}图是\Term{cartesian-diagram}{笛卡尔图}，且全部\Term{arrow}{箭头}都是\Term{injectivity}{单射}：
\[
\begin{tikzcd}
 H^1(S,P)\arrow[r,hook]&H^1(S,G)\\
 H^1(S,P\cap Q)\arrow[u,hook]\arrow[r,hook]&H^1(S,Q)\arrow[u,hook].
\end{tikzcd}
\]
\end{SGAresult}
证明（i）。由 2.3，映射 $H^1(S,L)\to H^1(S,P)$ 为\Term{bijection}{双射}；映射 $H^1(S,P)\to H^1_{t(P)}(S,G)$ \Term{surjectivity}{满射}（3.21），现证明它为\Term{injectivity}{单射}，即\Term{canonical}{典范}映射 $H^1(S,P)\to H^1(S,G)$ 为\Term{injectivity}{单射}。设 $Q$ 是 $P$ 下的\Term{torsor}{主丛}，$Q_1$ 是相伴的 $G$ 下\Term{torsor}{主丛}，$P'$ 和 $G'$ 是相应的 $P$ 和 $G$ 的\Term{twisting}{扭曲}\Term{form}{形式}。显然 $G'$ 是\Term{reductive-group}{约化 $S$ 群}，$P'$ 是它的\Term{parabolic-subgroup}{抛物子群}。$H^1(S,P)$ 中与 $Q$ 的类在 $H^1(S,G)$ 中有相同像的元素所成\Term{set}{集合}，自然等同于\Term{canonical}{典范}映射 $H^1(S,P')\to H^1(S,G')$ 的核；由\Term{cohomology}{上同调}\Term{exact-sequence}{正合序列}，后者又等同于
\SourcePage{318}
$G'(S)$ 在 $(G'/P')(S)$ 中的\Term{orbit}{轨道}集（这些\Term{nonabelian-cohomology}{非阿贝尔上同调}论证见 Giraud 的学位论文\footnote{编者注：见 [Gi71]，\S III.3。}）。但由 5.2，$G'(S)$ 在 $(G'/P')(S)$ 上\Term{transitive-action}{传递作用}。
证明（ii）：设 $Q$ 是 $G$ 下的\Term{torsor}{主齐性丛}，$G_Q$ 是 $G$ 相应的\Term{twisting}{扭曲}\Term{form}{形式}。根据定义（3.21），须证明：$G_Q$ 有 $t\cap t'$ 型\Term{parabolic-subgroup}{抛物子群}，当且仅当它有 $t$ 型及 $t'$ 型的\Term{parabolic-subgroup}{抛物子群}；这正是 3.8 和 5.7（i）的合取。最后，（iii）立即由（i）和（ii）得到。
现在列出 5.3 的一个推论。
\begin{SGAresult}{推论}{5.11}
设 $S$ 是\Term{semilocal-scheme}{半局部概形}，$G$ 是\Term{reductive-group}{约化 $S$ 群}，$P$ 是 $G$ 的\Term{parabolic-subgroup}{抛物子群}。若 $P\neq G$，则 $G$ 至少有三个彼此不同而与 $P$ 同型的\Term{parabolic-subgroup}{抛物子群}；换言之，$P\neq G$ 蕴含 $(G(S):P(S))\geq3$。
\end{SGAresult}
事实上，取 $G$ 中与 $P$ 相反的\Term{parabolic-subgroup}{抛物子群} $P'$（4.3.5（i））。因为 $P\neq G$，有 $\operatorname{rad}^{u}(P')\neq e$（例如由 4.3.2）。由 2.1，$\operatorname{rad}^{u}(P')(S)\neq e$；故可取 $u\in\operatorname{rad}^{u}(P')(S)$、$u\neq e$。于是 $\operatorname{int}(u)P\neq P$；由 5.3，存在与 $P$ 同型且与 $P$ 及 $\operatorname{int}(u)P$ 都相反的 $P_1$。于是 $P_1,P,\operatorname{int}(u)P$ 是 $G$ 的三个彼此不同、与 $P$ 同型的\Term{parabolic-subgroup}{抛物子群}。

\SGAsection{6}{抛物子群与分裂环面}
\begin{SGAresult}{命题}{6.1}
设 $S$ 是\Term{scheme}{概形}，$G$ 是\Term{reductive-group}{约化 $S$ 群}，$\mathfrak g=\operatorname{Lie}(G/S)$，$Q$ 是 $G$ 的\Term{split-torus}{分裂子环面}。写 $Q=\mathrm D_S(M)$，并设
\[
 \mathfrak g=\bigoplus_{\alpha\in M}\mathfrak g^\alpha
\]
是 $\mathfrak g$ 在 $Q$ 的\Term{action}{作用}下的\Term{resolution}{分解}。设 $M_1$ 是 $M$ 的一个\Term{subset}{子集}，满足 $0\in M_1$，且 $\alpha,\beta\in M_1\Rightarrow\alpha+\beta\in M_1$。

（i）存在唯一的 $G$ 的\Term{smooth}{光滑}\Term{subgroup}{子群} $H_{M_1}$，其\Term{fiber}{纤维}\Term{connected}{连通}，包含 $\operatorname{Centr}_G(Q)$，且其\Term{lie-algebra}{Lie 代数}为 $\bigoplus_{\alpha\in M_1}\mathfrak g^\alpha$。

（ii）有以下蕴含：
\begin{align*}
 M_1=\{0\}&\Longrightarrow H_{M_1}=\operatorname{Centr}_G(Q),\\
 M_1=-M_1&\Longrightarrow H_{M_1}\text{ 是约化群},\\
 M_1\cup(-M_1)=M&\Longrightarrow H_{M_1}\text{ 与 }H_{-M_1}\text{ 是 }G\text{ 的相反抛物子群，}\\
 &\hspace{2em}\text{具有公共 Levi 子群 }H_{M_1\cap-M_1}.
\end{align*}
\end{SGAresult}
为证明（i）和（ii），它们对（fpqc）\Term{topology}{拓扑}都是\Term{local-scheme}{局部}的，可假设 $Q$ 包含在 $G$ 的一个\Term{maximal-torus}{极大环面} $T$ 中；还可将 $G$ 相对于 $T$ \Term{split-group}{分裂}。断言（i）于是立即由第 XXII 讲 5.3.5、5.4.5 和 5.4.7 得到；（ii）的各断言由第 XXII 讲 5.3.5、5.10.1、5.11.3 及本讲 1.4 和 4.3.2 得到。

\SourcePage{319}
\begin{SGAresult}{推论}{6.2}
设 $S$ 是\Term{scheme}{概形}，$G$ 是\Term{reductive-group}{约化 $S$ 群}，$Q$ 是 $G$ 的\Term{split-torus}{分裂子环面}。存在 $G$ 的一个\Term{parabolic-subgroup}{抛物子群}，以 $\operatorname{Centr}_G(Q)$ 为\Term{levi-subgroup}{Levi 子群}。
\end{SGAresult}
事实上，写 $Q=\mathrm D_S(M)$，在群 $M$ 上选取一个全序结构，以 $M_1$ 表示 $M$ 中正元素的\Term{set}{集合}；群 $H_{M_1}$ 满足要求。
\begin{SGAresult}{推论}{6.3}
若\Term{reductive-group}{约化 $S$ 群} $G$ 有一个非\Term{center}{中心}的\Term{split-torus}{分裂子环面}，则它有一个真\Term{parabolic-subgroup}{抛物子群}（即不等于 $G$）。
\end{SGAresult}
由 5.9 和 5.10，从 6.2 得到：
\begin{SGAresult}{推论}{6.4}
设 $S$ 是\Term{scheme}{概形}，$G$ 是\Term{reductive-group}{约化 $S$ 群}，$Q$ 是 $G$ 的\Term{split-torus}{分裂子环面}。\Term{canonical}{典范}\Term{morphism}{态射} $G\to G/\operatorname{Centr}_G(Q)$ 是\Term{locally-trivial-bundle}{局部平凡纤维化}。
若 $S$ 半\Term{local-scheme}{局部}，则映射 $G(S)\to(G/\operatorname{Centr}_G(Q))(S)$ \Term{surjectivity}{满射}，映射 $H^1(S,\operatorname{Centr}_G(Q))\to H^1(S,G)$ \Term{injectivity}{单射}。
\end{SGAresult}

\textbf{6.5。}假设 $S$ \Term{connected}{连通}。若 $T$ 是一个 $S$ \Term{torus}{环面}，$T'$ 和 $T''$ 是 $T$ 的两个\Term{split-torus}{分裂子环面}，则其乘积 $T'\cdot T''$\footnote{编者注：采用了乘法记号，即将“它们的和 $T'+T''$”换成“它们的乘积 $T'\cdot T''$”。}也是 $T$ 的\Term{split-torus}{分裂子环面}。事实上，它等同于 $T'\times_ST''$ 对 $T'\cap T''$ 的商，而这个商由第 IX 讲 2.11 是\Term{split-torus}{分裂的}。因此 $T$ 有一个最大的\Term{split-torus}{分裂子环面}，记为 $T_{\mathrm{d\text{é}p}}$。

\begin{SGAresult}{引理}{6.6}
设 $S$ 是\Term{connected}{连通}\Term{scheme}{概形}，$T$ 是\Term{isotrivial}{等平凡（isotrivial）的} $S$ \Term{torus}{环面}，$T_{\mathrm{d\text{é}p}}$ 是其最大的\Term{split-torus}{分裂子环面}。以下条件\Term{equivalence}{等价}：
（i）存在不同于 $e$ 的\Term{homomorphism}{同态} $T\to\mathbb G_{m,S}$。
（ii）$T_{\mathrm{d\text{é}p}}\neq e$。
\end{SGAresult}
因为假设 $T$ \Term{isotrivial}{等平凡}，存在\Term{finite-group}{有限群} $\Gamma$、群为 $\Gamma$ 的\Term{connected}{连通} Galois 主\Term{covering}{覆盖} $S'\to S$ 及\Term{isomorphism}{同构} $T_{S'}\simeq\mathrm D_{S'}(M)$；于是 $M$ 配有 $\Gamma$ 模结构，有自然\Term{isomorphism}{同构} $\operatorname{Hom}_S(T,\mathbb G_{m,S})=H^0(\Gamma,M)$。
另一方面，设 $V$ 是 $M\otimes\mathbb Q$ 中由形如 $g(m)-m$（$g\in\Gamma$、$m\in M$）的元素张成的向量子\Term{space}{空间}。立即可验 $(T_{\mathrm{d\text{é}p}})_{S'}$ 等同于 $\mathrm D_{S'}(M/(M\cap V))$。因此断言（i）\Term{equivalence}{等价}于 $H^0(\Gamma,M)\neq0$，也\Term{equivalence}{等价}于 $H^0(\Gamma,M\otimes\mathbb Q)\neq0$；断言（ii）\Term{equivalence}{等价}于 $M\neq M\cap V$，也\Term{equivalence}{等价}于 $M\otimes\mathbb Q\neq V$。但有
\[
 M\otimes\mathbb Q=H^0(\Gamma,M\otimes\mathbb Q)\oplus V,
\]
这是立即可验的（考虑\Term{projection-operator}{投影算子} $M\otimes\mathbb Q\to H^0(\Gamma,M\otimes\mathbb Q)$，它把 $x$ 映到 $x$ 在 $\Gamma$ 下各变换的平均值）。
\begin{SGAresult}{引理}{6.7}
设 $S$ 是\Term{scheme}{概形}，$G$ 是\Term{reductive-group}{约化 $S$ 群}，$P$ 是 $G$ 的\Term{parabolic-subgroup}{抛物子群}且 $P\neq G$，$L$ 是 $P$ 的一个\Term{levi-subgroup}{Levi 子群}，$Q$ 是其\Term{radical}{根基}。存在不同于 $e$ 的\Term{homomorphism}{同态} $Q\to\mathbb G_{m,S}$。
\end{SGAresult}
\SourcePage{320}
考虑 $P$ 的\Term{unipotent-radical}{幺幂根基} $U$；它在 $\operatorname{int}(P)$ 下不变，故在 $\operatorname{int}(Q)$ 下不变。考虑\Term{invertible}{可逆} $\mathscr O_S$ 模 $\det(\operatorname{Lie}(U))$，即\Term{locally-free}{局部自由} $\mathscr O_S$ 模 $\operatorname{Lie}(U)$ 的“最高\Term{exterior-power}{外幂}”。\Term{adjoint-representation}{伴随表示}定义了\Term{group-homomorphism}{群同态}
\[
 f:Q\longrightarrow\operatorname{Aut}(\det(\operatorname{Lie}(U)))=\mathbb G_{m,S}.
\]
若 $P\neq G$，则 $U\neq e$。选取 $s\in S$，使 $U_s\neq e$。将 $G_s$ 相对于一个包含 $Q_s$ 的\Term{maximal-torus}{极大环面}\Term{split-group}{分裂}，立即看出 $f_s\neq e$。

\begin{SGAresult}{命题}{6.8}
设 $S$ 是\Term{connected}{连通}\Term{semilocal-scheme}{半局部概形}，$G$ 是\Term{reductive-group}{约化 $S$ 群}，$P$ 是 $G$ 的\Term{parabolic-subgroup}{抛物子群}，$L$ 是 $P$ 的\Term{levi-subgroup}{Levi 子群}，$Q$ 是其\Term{radical}{根基}，$Q_{\mathrm{d\text{é}p}}$ 是 $Q$ 的最大\Term{split-torus}{分裂子环面}（即 $L$ 的最大\Term{center}{中心}\Term{split-torus}{分裂子环面}）。则
\[
 L=\operatorname{Centr}_G(Q_{\mathrm{d\text{é}p}}).
\]
\end{SGAresult}
置 $L'=\operatorname{Centr}_G(Q_{\mathrm{d\text{é}p}})$；它是 $G$ 中包含 $L$ 的\Term{reductive-group}{约化子群}；而且 $P'=P\cap L'$ 是 $L'$ 的\Term{parabolic-subgroup}{抛物子群}，以 $L$ 为\Term{levi-subgroup}{Levi 子群}（1.20）。若 $L'\neq L$，则 $L'\not\subset P$（因 $L$ 是 $P$ 的极大\Term{reductive-group}{约化子群}，见 1.7），故 $P'\neq L'$。
设 $G_1$ 是 $G$ 的\Term{derived-group}{导群}，$P_1=P'\cap G_1$。由 1.19，$P_1$ 是\Term{semisimple-group}{半单群} $G_1$ 的\Term{parabolic-subgroup}{抛物子群}，$L_1=L\cap G_1$ 是它的\Term{levi-subgroup}{Levi 子群}，并且
\[
 Q_1=\operatorname{rad}(L_1)=(\operatorname{rad}(L)\cap G_1)^0=(Q\cap G_1)^0.
\]
因为 $L_1$ 有一个\Term{maximal-torus}{极大环面} $T_1$（第 XIV 讲 3.20），且它\Term{isotrivial}{等平凡}（第 XXIV 讲 4.1.5），$T_1$ 的\Term{subtorus}{子环面} $Q_1$ 也\Term{isotrivial}{等平凡}（第 IX 讲 2.11）。因 $P_1\neq G_1$，可应用 6.7 和 6.6，故 $(Q_1)_{\mathrm{d\text{é}p}}\neq e$，从而 $(Q_{\mathrm{d\text{é}p}}\cap G_1)^0\neq e$，于是 $Q_{\mathrm{d\text{é}p}}\not\subset\operatorname{rad}(L')$（因为 $\operatorname{rad}(L')\cap G_1$ \Term{finite}{有限}），这与 $L'$ 的定义矛盾。

\begin{SGAresult}{推论}{6.9}
设 $S$ 是\Term{connected}{连通}\Term{semilocal-scheme}{半局部概形}，$G$ 是\Term{reductive-group}{约化 $S$ 群}，$Q$ 是 $G$ 的\Term{critical-torus}{临界子环面}（即 $\operatorname{rad}(\operatorname{Centr}_G(Q))=Q$）。$\operatorname{Centr}_G(Q)$ 是 $G$ 的某个\Term{parabolic-subgroup}{抛物子群}的\Term{levi-subgroup}{Levi 子群}，当且仅当 $\operatorname{Centr}_G(Q)=\operatorname{Centr}_G(Q_{\mathrm{d\text{é}p}})$，即 $\operatorname{Lie}(G)^Q=\operatorname{Lie}(G)^{Q_{\mathrm{d\text{é}p}}}$。
\end{SGAresult}
这由 6.2 和 6.8 得到。
\begin{SGAresult}{推论}{6.10}
设 $S$ 是\Term{connected}{连通}\Term{semilocal-scheme}{半局部概形}，$G$ 是\Term{reductive-group}{约化 $S$ 群}，$L$ 是 $G$ 的\Term{subgroup}{子群}。以下条件\Term{equivalence}{等价}：
（i）存在 $G$ 的\Term{parabolic-subgroup}{抛物子群}，以 $L$ 为\Term{levi-subgroup}{Levi 子群}。
（ii）存在 $G$ 的\Term{split-torus}{分裂子环面}，以 $L$ 为\Term{centralizer}{中心化子}。
（iii）存在\Term{homomorphism}{同态} $\mathbb G_{m,S}\to G$，以 $L$ 为\Term{centralizer}{中心化子}。
\end{SGAresult}
由 6.8 有（i）$\Rightarrow$（ii），由 6.1 有（iii）$\Rightarrow$（i）；还须证明（ii）$\Rightarrow$（iii）。因此假设 $L=\operatorname{Centr}_G(Q)$，其中 $Q=\mathrm D_S(M)$。写 $\mathfrak g=\operatorname{Lie}(G/S)$ 及
\[
 \mathfrak g=\bigoplus_{\alpha\in M}\mathfrak g^\alpha,
\]
设 $R$ 是使 $\mathfrak g^\alpha\neq0$ 的 $\alpha\in M-\{0\}$ 所成的\Term{set}{集合}。
\SourcePage{321}
因为 $R$ \Term{finite}{有限}且不含 $0$，存在\Term{homomorphism}{同态} $u:M\to\mathbb Z$，使每个 $\alpha\in R$ 都满足 $u(\alpha)\neq0$。由\Term{duality}{对偶性}，$u$ 给出\Term{homomorphism}{同态} $\mathbb G_{m,S}\to Q$，故给出\Term{homomorphism}{同态} $f:\mathbb G_{m,S}\to G$。有 $\operatorname{Centr}_G(f)\supset\operatorname{Centr}_G(Q)$；它们是 $G$ 的两个\Term{fiber}{纤维}\Term{connected}{连通}的\Term{smooth}{光滑}\Term{subgroup}{子群}；它们的\Term{lie-algebra}{Lie 代数}相同（都等于 $\mathfrak g^0$）；由通常的论证，它们相等。

\begin{SGAresult}{推论}{6.11}
设 $S$ 是\Term{connected}{连通}\Term{semilocal-scheme}{半局部概形}，$G$ 是\Term{reductive-group}{约化 $S$ 群}。映射
\[
 L\longmapsto\operatorname{rad}(L)_{\mathrm{d\text{é}p}},\qquad
 Q\longmapsto\operatorname{Centr}_G(Q)
\]
在如下两个\Term{set}{集合}之间给出互逆的、反转自然序结构的\Term{bijection}{双射}：一方是 $G$ 中作为 $G$ 的\Term{parabolic-subgroup}{抛物子群}的\Term{levi-subgroup}{Levi 子群}的\Term{subgroup}{子群} $L$ 所成的\Term{set}{集合}；另一方是 $G$ 中满足 $\operatorname{rad}(\operatorname{Centr}_G(Q))_{\mathrm{d\text{é}p}}=Q$ 的\Term{split-torus}{分裂子环面} $Q$ 所成的\Term{set}{集合}。
\end{SGAresult}
\begin{SGAresult}{推论}{6.12}
设 $S$ 是\Term{connected}{连通}\Term{semilocal-scheme}{半局部概形}，$G$ 是\Term{reductive-group}{约化 $S$ 群}。考虑以下断言：
（i）存在 $G$ 的一个不同于 $G$ 的\Term{parabolic-subgroup}{抛物子群}。
（ii）$G$ 有一个非\Term{center}{中心}的\Term{split-torus}{分裂子环面}。
（ii bis）$G$ 有一个\Term{relative-dimension}{相对维数}为 $1$ 的非\Term{center}{中心}\Term{split-torus}{分裂子环面}。
（iii）存在\Term{group-homomorphism}{群同态} $\mathbb G_{a,S}\to G$，它是\Term{closed-immersion}{闭浸入}。
则有（i）$\Leftrightarrow$（ii）$\Leftrightarrow$（ii bis）$\Rightarrow$（iii）。
\end{SGAresult}
唯一的新断言是（i）$\Rightarrow$（iii）。设 $P$ 是 $G$ 中不同于 $G$ 的\Term{parabolic-subgroup}{抛物子群}。则 $U=\operatorname{rad}^{u}(P)\neq e$。考虑 $U$ 的合成列（2.1）中最后一个非平凡\Term{subgroup}{子群} $U_n$。有\Term{isomorphism}{同构} $U_n\simeq\mathbf W(E_n)$，其中 $E_n$ 是\Term{locally-free}{局部自由} $\mathscr O_S$ 模，故为\Term{free-module}{自由模}\footnote{编者注：因为假设 $S$ 半\Term{local-scheme}{局部}且\Term{connected}{连通}。}。因为 $E_n\neq0$，存在\Term{local-scheme}{局部}为\Term{direct-summand}{直和因子}的\Term{monomorphism}{单态射} $\mathscr O_S\to E_n$，故有\Term{closed-immersion}{闭浸入}
$\mathbb G_{a,S}=\mathbf W(\mathscr O_S)\hookrightarrow\mathbf W(E_n)\simeq U_n$，立即推出（iii）。

\textbf{注记 6.12.1。}当 $S$ 是\Term{characteristic}{特征} $0$ 的域的谱时，由 Jacobson--Morozov 定理，（iii）$\Rightarrow$（ii bis）。这时上述四个条件\Term{equivalence}{等价}（“Godement 判据”，见 [BT65]，8.5）。\footnote{原注（$*$）：更一般地，当 $S$ 是\Term{perfect-field}{完美域}的谱时也成立（Tits）。编者注（30）：这是文章 [BT71] 的推论 3.7。}
\begin{SGAresult}{定义}{6.13}
设 $S$ 是\Term{connected}{连通}\Term{semilocal-scheme}{半局部概形}，$G$ 是\Term{reductive-group}{约化 $S$ 群}。若 $G$ 不含任何不等于 $e$ 的\Term{split-torus}{分裂子环面}，就称 $G$ \Term{anisotropic-group}{各向异性}。
\end{SGAresult}
\begin{SGAresult}{推论}{6.14}
设 $S$ 是\Term{connected}{连通}\Term{semilocal-scheme}{半局部概形}。\Term{reductive-group}{约化 $S$ 群} $G$ \Term{anisotropic-group}{各向异性}，当且仅当它没有任何 $P\neq G$ 的\Term{parabolic-subgroup}{抛物子群}，且其\Term{radical}{根基}\Term{anisotropic-group}{各向异性}。
\end{SGAresult}
利用 6.6、第 XXIV 讲 4.1.5 和第 XXII 讲 6.2，得到：
\SourcePage{322}
\begin{SGAresult}{推论}{6.15}
设 $S$ 是\Term{connected}{连通}\Term{semilocal-scheme}{半局部概形}，$G$ 是\Term{isotrivial}{等平凡的}\Term{reductive-group}{约化 $S$ 群}（例如 $G$ \Term{semisimple}{半单}，或 $S$ \Term{normal}{正规}（第 X 讲 5.16））。$G$ \Term{anisotropic-group}{各向异性}，当且仅当 $G$ 没有 $P\neq G$ 的\Term{parabolic-subgroup}{抛物子群}，且 $\operatorname{Hom}_{S\text{-gr}}(G,\mathbb G_{m,S})=e$。
\end{SGAresult}
\begin{SGAresult}{命题}{6.16}
设 $S$ 是\Term{connected}{连通}\Term{semilocal-scheme}{半局部概形}，$G$ 是\Term{reductive-group}{约化 $S$ 群}。$G$ 的极大\Term{split-torus}{分裂子环面}正是 $G$ 的\Term{minimal-parabolic}{极小抛物子群}的\Term{levi-subgroup}{Levi 子群}的最大\Term{center}{中心}\Term{split-torus}{分裂子环面}。任意两个这样的\Term{torus}{环面}由 $G(S)$ 的一个元素\Term{conjugacy}{共轭}。
\end{SGAresult}
设 $Q$ 是 $G$ 的极大\Term{split-torus}{分裂子环面}。\footnote{编者注：已将下面一句展开说明。}由 6.2，$L=\operatorname{Centr}_G(Q)$ 是 $G$ 的某个\Term{parabolic-subgroup}{抛物子群} $P$ 的\Term{levi-subgroup}{Levi 子群}；因 $Q\subset\operatorname{rad}(\operatorname{Centr}_G(Q))_{\mathrm{d\text{é}p}}$，$Q$ 的极大性推出 $Q=\operatorname{rad}(\operatorname{Centr}_G(Q))_{\mathrm{d\text{é}p}}$。由 6.11，$L$ 是 $G$ 的\Term{parabolic-subgroup}{抛物子群}的\Term{levi-subgroup}{Levi 子群}所成\Term{set}{集合}的极小元素，故由 1.20，$P$ 是 $G$ 的\Term{minimal-parabolic}{极小抛物子群}。由 5.7 和 5.5（iv），任意两个这样的\Term{subtorus}{子环面} $Q$ 由 $G(S)$ 的一个\Term{section}{截面}\Term{conjugacy}{共轭}。$Q$ 以及对子 $(P,L)$ 的\Term{conjugacy}{共轭}性随即推出 6.16 的第一个断言。
\begin{SGAresult}{推论}{6.17}
设 $S$ 是\Term{connected}{连通}\Term{semilocal-scheme}{半局部概形}，$P$ 和 $P'$ 是两个处于\Term{standard-position}{标准位置}的\Term{minimal-parabolic}{极小抛物子群}（4.5.1.1）。则 $P\cap P'$ 包含 $P$ 和 $P'$ 的一个公共\Term{levi-subgroup}{Levi 子群}。
\end{SGAresult}
事实上，由 4.5.1（v），$P\cap P'$ 包含 $G$ 的一个\Term{maximal-torus}{极大环面} $T$；设 $L$ 是 $P$ 中唯一包含 $T$ 的\Term{levi-subgroup}{Levi 子群}。由 1.21，有
\[
 \operatorname{rad}(P)\cap T=\operatorname{rad}(P)\cap L=\operatorname{rad}(L),
\]
故 $\operatorname{rad}(P)\cap T$ 包含 $\operatorname{rad}(L)_{\mathrm{d\text{é}p}}$；后者是 $G$ 的极大\Term{split-torus}{分裂子环面}，故必等于 $T_{\mathrm{d\text{é}p}}$。所以 $L=\operatorname{Centr}_G(T_{\mathrm{d\text{é}p}})$，由对称性，$L$ 也是 $P'$ 的\Term{levi-subgroup}{Levi 子群}。

\textbf{注记 6.18。}由 1.21，$G$ 的\Term{parabolic-subgroup}{抛物子群} $P$ 是\Term{minimal-parabolic}{极小的}，当且仅当 $\operatorname{rad}(P)$ 包含 $G$ 的一个极大\Term{split-torus}{分裂子环面}；这时由 6.17，若 $T$ 是 $P$ 的一个\Term{maximal-torus}{极大环面}，则 $T_{\mathrm{d\text{é}p}}$ 是 $G$ 和 $\operatorname{rad}(P)$ 的一个\Term{maximal-torus}{极大环面}，且 $\operatorname{Centr}_G(T_{\mathrm{d\text{é}p}})$ 是 $P$ 的\Term{levi-subgroup}{Levi 子群}。此外，$P$ 的每个\Term{levi-subgroup}{Levi 子群}都以这种方式得到。

\SGAsection{7}{相对根数据}
本节中，$S$ 表示\Term{nonempty}{非空}\Term{connected}{连通}\Term{semilocal-scheme}{半局部概形}，$G$ 是\Term{reductive-group}{约化 $S$ 群}，$Q$ 是 $G$ 的极大\Term{split-torus}{分裂子环面}，$L$ 是 $Q$ 在 $G$ 中的\Term{centralizer}{中心化子}，即 $L=\operatorname{Centr}_G(Q)$。
\SourcePage{323}
\textbf{7.1。}因为 $Q$ 是 $L$ 的最大\Term{center}{中心}平凡\Term{subtorus}{子环面}，$G(S)$ 中\Term{normalization}{正规化} $L$ 的每个\Term{section}{截面}也\Term{normalization}{正规化} $Q$。因此（见 7.1.1）
\[
 \operatorname{Norm}_G(L)(S)=\operatorname{Norm}_G(Q)(S).
\]
另一方面，6.4 已说明映射 $G(S)\to(G/L)(S)$ \Term{surjectivity}{满射}。于是有\Term{canonical}{典范}等同
\[
 W_G(Q)(S)=(\operatorname{Norm}_G(Q)/\operatorname{Centr}_G(Q))(S)
 \simeq\operatorname{Norm}_G(L)(S)/L(S).
\]
用 $M$ 表示群 $\operatorname{Hom}_{S\text{-gr}}(Q,\mathbb G_{m,S})$，于是有\Term{canonical-isomorphism}{典范同构} $Q\simeq\mathrm D_S(M)$。以 $W$ 表示由 $W_G(Q)(S)$ 定义的 $M$ 的\Term{automorphism-group}{自同构群}。因此有\Term{isomorphism}{同构}
\[
 W\simeq W_G(Q)(S)\simeq\operatorname{Norm}_G(L)(S)/L(S).
\]

\textbf{7.1.1。}一般并无 $\operatorname{Norm}_G(L)=\operatorname{Norm}_G(Q)$。例如，取 $S$ 为一个有二次\Term{group-extension}{扩张} $k'$ 的域 $k$ 的谱，取 $G$ 为酉群 $G^{q\mathrm{\text{É}p}}_{k'/k}(A_2)$（见第 XXIV 讲 3.11.2）。
因 $G$ 的\Term{minimal-parabolic}{极小抛物子群}就是它的\Term{borel-subgroup}{Borel 子群}，其\Term{levi-subgroup}{Levi 子群}都是\Term{maximal-torus}{极大环面}，且有 $(\operatorname{Norm}_G(L)/L)(\bar k)=\mathfrak S_3$。
另一方面，因为 $G$ 不\Term{split-group}{分裂}，$G$ 的极大\Term{split-torus}{分裂环面}\Term{dimension}{维数} $\leq1$，故\Term{isomorphism}{同构}于 $\mathbb G_{m,k}$。因 $\operatorname{Norm}_G(Q)/L$ 在 $Q$ 上\Term{faithful-action}{忠实作用}，有 $(\operatorname{Norm}_G(Q)/Q)(\bar k)=\mathbb Z/2\mathbb Z$。

\textbf{7.2。}若 $P$ 是 $G$ 中以 $L$ 为\Term{levi-subgroup}{Levi 子群}的\Term{parabolic-subgroup}{抛物子群}（由 6.2 存在），则 $P$ 必为\Term{minimal-parabolic}{极小的}（见 6.18）。由 $G$ 的\Term{minimal-parabolic}{极小抛物子群}的\Term{conjugacy}{共轭}性（5.7）、一个\Term{parabolic-subgroup}{抛物子群}的\Term{levi-subgroup}{Levi 子群}的\Term{conjugacy}{共轭}性（1.8），以及等式 $P=\operatorname{Norm}_G(P)$ 和 $\operatorname{Norm}_G(L)\cap P=L$（1.6），得到：$G$ 中以 $L$ 为\Term{levi-subgroup}{Levi 子群}的（\Term{minimal-parabolic}{极小}）\Term{parabolic-subgroup}{抛物子群}所成的\Term{set}{集合}，是群 $W$ 下的\Term{principal-homogeneous-space}{主齐性空间}。

\textbf{7.3。}$G$ 的\Term{lie-algebra}{Lie 代数}在 $Q$ 的\Term{action}{作用}下\Term{resolution}{分解}为
\[
 \operatorname{Lie}(G)=\operatorname{Lie}(L)\oplus\bigoplus_{\alpha\in R}\operatorname{Lie}(G)^\alpha,
\]
其中 $R$ 是使 $\operatorname{Lie}(G)^\alpha\neq0$ 的 $Q$ 的非零\Term{character}{特征标}所成的\Term{set}{集合}（$G$ 相对于 $Q$ 的\Term{root}{根}）。
以 $M^*$ 表示群 $\operatorname{Hom}_{S\text{-gr}}(\mathbb G_{m,S},Q)$；它与 $M$ \Term{duality}{对偶}，$W$ 通过\Term{transport-of-structure}{结构搬运}在其上自然\Term{action}{作用}。
\begin{SGAresult}{定理}{7.4}
沿用 7.3 的记号，存在唯一的映射 $\alpha\mapsto\alpha^*$，从 $R$ 到 $M^*$，在 $(M,M^*)$ 中定义一个\Term{weyl-group}{Weyl 群}为 $W$ 的\Term{root-datum}{根数据}（第 XXI 讲 1.1）。
此外，$G$ 中以 $L$ 为\Term{levi-subgroup}{Levi 子群}的\Term{parabolic-subgroup}{抛物子群} $P$ 与 $R$ 的\Term{positive-roots}{正根系} $R_+$ 由关系
\[
 \operatorname{Lie}(P)=\operatorname{Lie}(L)\oplus\bigoplus_{\alpha\in R_+}\operatorname{Lie}(G)^\alpha
\]
一一对应。
\end{SGAresult}
\SourcePage{324}
\textbf{7.4.1。}先假设所求映射 $\alpha\mapsto\alpha^*$ 的存在性已获证明。由第 XXI 讲 3.4.10，$s_\alpha$ 是 $W$ 中唯一满足“对每个 $m\in M$，$s_\alpha(m)-m$ 为 $\alpha$ 的有理倍数”的元素，这说明 $s_\alpha$ 由 $\alpha$ 决定。因这时有\footnote{编者注：已将 $\alpha^*(m)$ 改为 $\alpha^*(m)\alpha$。}$\alpha^*(m)\alpha=m-s_\alpha(m)$，可见 $\alpha^*$ 由 $\alpha$ 决定，从而证明了映射 $\alpha\mapsto\alpha^*$ 的唯一性。

\textbf{7.4.2。}设 $\alpha\in R$，以 $L_\alpha$（分别以 $H_\alpha$、$H_{-\alpha}$）表示 $G$ 中唯一包含 $L$、\Term{smooth}{光滑}且\Term{fiber}{纤维}\Term{connected}{连通}、并满足下式的\Term{subgroup}{子群}（见 6.1（i））：
\begin{align*}
 \operatorname{Lie}(L_\alpha)&=\operatorname{Lie}(L)\oplus\bigoplus_{\gamma\in\mathbb Z\alpha\cap R}\operatorname{Lie}(G)^\gamma,\\
 \text{分别为}\quad\operatorname{Lie}(H_\alpha)&=\operatorname{Lie}(L)\oplus\bigoplus_{\gamma\in\mathbb N\alpha\cap R}\operatorname{Lie}(G)^\gamma,\\
 \text{分别为}\quad\operatorname{Lie}(H_{-\alpha})&=\operatorname{Lie}(L)\oplus\bigoplus_{\gamma\in-\mathbb N\alpha\cap R}\operatorname{Lie}(G)^\gamma.
\end{align*}
$L_\alpha$ 是 $G$ 的\Term{reductive-group}{约化子群}，$H_\alpha$ 和 $H_{-\alpha}$ 是它的\Term{parabolic-subgroup}{抛物子群}，以 $L$ 为\Term{levi-subgroup}{Levi 子群}，且 $H_\alpha$ 与 $H_{-\alpha}$ 相对于 $L$ 相反（见 6.1（ii））。
因此由 7.2，存在 $s_\alpha\in\operatorname{Norm}_{L_\alpha}(Q)(S)/L(S)\subset W$，使 $s_\alpha(H_\alpha)=H_{-\alpha}$。有 $s_\alpha(\alpha)=-\alpha$（因为 $\alpha$，分别为 $-\alpha$，是出现在 $H_\alpha$，分别为 $H_{-\alpha}$，中的 $R$ 的元素的公因子），并有 $s_\alpha^2=\mathrm{id}$（因为 $s_\alpha^2(H_\alpha)$ 和 $H_\alpha$ 都相对于 $L$ 与 $H_{-\alpha}$ 相反）。于是构造了满足下列性质的 $s_\alpha\in W$：
\[
 s_\alpha(\alpha)=-\alpha,\qquad s_\alpha^2=\mathrm{id};\tag{x}
\]
\[
 s_\alpha\text{ 可由 }L_\alpha(S)\text{ 的一个元素代表。}\tag{xx}
\]
还须指出，$s_\alpha$ 由 $\alpha$ \Term{canonical}{典范}地构造而来，特别有
\[
 \text{对每个 }w\in W,\quad ws_\alpha w^{-1}=s_{w(\alpha)}.\tag{xxx}
\]

\textbf{7.4.3。}现在拟证明断言
\[
 \text{对每个 }m\in M,\quad s_\alpha(m)-m\in\mathbb Z\alpha.\tag{xxxx}
\]
因为 $S$ \Term{connected}{连通}，此断言对（fpqc）\Term{topology}{拓扑}是\Term{local-scheme}{局部}的。因此可以假设 $L_\alpha=G_1$ 相对于 $L$ 的一个\Term{maximal-torus}{极大环面} $T_1$ 可\Term{split-group}{分裂}。设 $(G_1,T_1,M_1,R_1)$ 是这样一个\Term{splitting}{分裂结构}。\Term{monomorphism}{单态射} $Q\to T_1$ 将 $M$ 等同于 $M_1$ 的一个商，以 $p:M_1\to M$ 表示\Term{canonical}{典范}映射。
$R_1$ 在 $p$ 下的像由 $G_1$ 相对于 $Q$ 的\Term{root}{根}（因而是 $R$ 中 $\alpha$ 的整数倍）及可能的零组成；因此 $p(R_1)\subset\mathbb Z\alpha$。由（xx），存在 $\operatorname{Norm}_{G_1}(L)(S)$ 的一个元素，在 $Q$ 上诱导 $s_\alpha$。由第 XXII 讲 5.10.10，因此存在\Term{section}{截面} $w\in(W_1)_S(S)$，在 $Q$ 上诱导 $s_\alpha$（以 $W_1$ 表示\Term{root-datum}{根数据} $(M_1,R_1,\ldots)$ 的\Term{weyl-group}{Weyl 群}）。
\SourcePage{325}
适当\Term{restriction}{限制} $S$，可假设存在 $w\in W_1$，在 $Q$ 上诱导 $s_\alpha$，故对每个 $m_1\in M_1$ 满足 $p(w(m_1))=s_\alpha(p(m_1))$。但根据 $W_1$ 的定义，$w$ 是关于 $R_1$ 的元素的若干反射的乘积，故 $w(m_1)-m_1$ 是 $R_1$ 的元素的整系数线性组合。因此 $s_\alpha(p(m_1))-p(m_1)$ 是 $p(R_1)\subset\mathbb Z\alpha$ 的元素的整系数线性组合，故是 $\alpha$ 的整数倍，所求得证。

\textbf{7.4.4。}因此可由\footnote{编者注：已将 $\alpha^*(m)$ 改为 $\alpha^*(m)\alpha$。}
\[
 \alpha^*(m)\alpha=m-s_\alpha(m)
\]
定义元素 $\alpha^*\in M^*$。由（x），有 $(\alpha^*,\alpha)=2$；另一方面，由（xxx），对每对 $(\alpha,\alpha')\in R\times R$，有 $s_\alpha(\alpha')\in R$ 及 $s_\alpha(\alpha'^*)=s_\alpha(\alpha')^*$。这证明所构造的映射 $\alpha\mapsto\alpha^*$ 确在 $(M,M^*)$ 中定义了\Term{root-datum}{根数据}（见第 XXI 讲 1.1）。

\textbf{7.4.5。}设 $W_0$ 是此\Term{root-datum}{根数据}的\Term{weyl-group}{Weyl 群}（由各 $s_\alpha$ 生成的 $M$ 的变换群）；有 $W_0\subset W$。
另一方面，设 $\geq$ 是自由\Term{abelian-group}{阿贝尔群} $M$ 上的一个全\Term{order-relation}{序关系}；置 $R_+=\{\alpha\in R\mid\alpha\geq0\}$。已知 $R_+$ 是 $R$ 的一个\Term{positive-roots}{正根系}。设 $w\in W$ 由 $n\in\operatorname{Norm}_G(L)(S)=\operatorname{Norm}_G(Q)(S)$ 代表。置 $P=H_{R_+}$（6.1 的记号）；由该处，$P$ 是 $G$ 的\Term{parabolic-subgroup}{抛物子群}，以 $L$ 为\Term{levi-subgroup}{Levi 子群}。显然 $\operatorname{int}(n)P=H_{w(R_+)}$。由 7.3，$w(R_+)=R_+$ 蕴含 $w=e$。因群 $W_0$ 在 $R$ 的\Term{positive-roots}{正根系}上\Term{transitive-action}{传递作用}（第 XXI 讲 3.3.7），且 $R_+$ 在 $W$ 中的\Term{stabilizer}{稳定子}为恒等元，立即推出 $W=W_0$。还推出 $W=W_0$ 既在 $R$ 的\Term{positive-roots}{正根系}的\Term{set}{集合}上、又在 $G$ 中以 $L$ 为\Term{levi-subgroup}{Levi 子群}的\Term{parabolic-subgroup}{抛物子群}的\Term{set}{集合}上简\Term{simply-transitive}{单传递作用}，这给出 7.4 的最后一个断言。证毕。

\textbf{7.5。}若 $P$ 和 $P_1$ 是 $G$ 的两个\Term{minimal-parabolic}{极小抛物子群}，分别以 $L$ 和 $L_1$ 为\Term{levi-subgroup}{Levi 子群}，以 $Q$ 和 $Q_1$ 表示 $L$ 和 $L_1$ 的极大\Term{center}{中心}\Term{split-torus}{分裂环面}，则对子 $(P,Q)$ 和 $(P_1,Q_1)$ \Term{conjugacy}{共轭}：存在 $g\in G(S)$，使 $\operatorname{int}(g)P=P_1$、$\operatorname{int}(g)Q=Q_1$。
事实上，$P$ 和 $P_1$ \Term{conjugacy}{共轭}（5.7），故可假设 $P=P_1$；这时 $L$ 和 $L_1$ 由 $P(S)$ 的一个\Term{section}{截面}\Term{conjugacy}{共轭}（1.8）。此外，若 $g$ 和 $g'$ 是 $G$ 中使对子 $(P,Q)$ 与 $(P_1,Q_1)$ \Term{conjugacy}{共轭}的两个\Term{section}{截面}，则 $g'^{-1}g$ \Term{normalization}{正规化} $P$ 和 $Q$，因而\Term{normalization}{正规化} $P$ 和 $L$；但
\[
 \operatorname{Norm}_G(P)\cap\operatorname{Norm}_G(L)
 =P\cap\operatorname{Norm}_G(L)=L=\operatorname{Centr}_G(Q).
\]
因此由 $\operatorname{int}(g)$ 诱导的\Term{isomorphism}{同构} $Q\xrightarrow{\sim}Q_1$ 与 $g$ 无关。
设 $\mathscr R$ 和 $\mathscr R_1$ 是由 7.4 在 $\operatorname{Hom}_{S\text{-gr}}(Q,\mathbb G_{m,S})$ 和 $\operatorname{Hom}_{S\text{-gr}}(Q_1,\mathbb G_{m,S})$ 中定义的\Term{root-datum}{根数据}，$R_+$ 和 $R_{1+}$ 是对应于 $P$ 和 $P_1$ 的\Term{positive-roots}{正根系}。上面定义的\Term{canonical-isomorphism}{典范同构} $Q\xrightarrow{\sim}Q_1$ 将
\SourcePage{326}
$(\mathscr R,R_+)$ 变为 $(\mathscr R_1,R_{1+})$。立即推出，可以利用上述\Term{isomorphism}{同构}的传递系将不同的 $(\mathscr R,R_+)$ 认同，从而定义 $G$ 在 $S$ 上的相对\Term{pinned-root-datum}{带钉扎根数据}\footnote{编者注：回顾（第 XXIII 讲 1.5），\Term{pinned-root-datum}{带钉扎根数据}是配有一个\Term{positive-roots}{正根系}（或\Term{simple-roots}{单根系}）的选择的\Term{root-datum}{根数据}。}。
从现在起，将这个\Term{pinned-root-datum}{带钉扎根数据}记为
$(\underline M,\underline M^*,\underline R,\underline R^*,\underline R_+)=\mathscr R(G/S)$；于是对每个如上的对子 $P\supset Q$，有\Term{canonical-isomorphism}{典范同构} $\underline M\xrightarrow{\sim}\operatorname{Hom}_{S\text{-gr}}(Q,\mathbb G_{m,S})$，将 $\underline R$（分别为 $\underline R_+$）变为 $G$（分别为 $P$）相对于 $Q$ 的\Term{root}{根}的\Term{set}{集合}，并将 $W(\underline R)$ 变为 $W_G(Q)(S)$。

\textbf{7.6。}仍设 $Q$ 是 $G$ 的极大\Term{split-torus}{分裂环面}，$P$ 是 $G$ 中以 $\operatorname{Centr}_G(Q)$ 为\Term{levi-subgroup}{Levi 子群}的（\Term{minimal-parabolic}{极小}）\Term{parabolic-subgroup}{抛物子群}，$(M,M^*,R,R^*,R_+)$ 是相应的\Term{pinned-root-datum}{带钉扎根数据}（7.4），$\Delta$ 是 $R_+$ 的\Term{simple-roots}{单根集}。对每个 $A\subset\Delta$，设 $R_A\subset R$ 为\Term{set}{集合}
\[
 R_A=R_+\cup(\mathbb ZA\cap R_-),
\]
它由所有\Term{positive-root}{正根}及可写成 $A$ 的元素的线性组合的\Term{negative-root}{负根}组成。这是\Term{root}{根}的一个闭集（第 XXI 讲 3.1.4），而每个包含 $R_+$ 的闭集都唯一地具有这种\Term{form}{形式}（第 XXI 讲 3.3.10）。由 6.1，存在唯一的 $G$ 的\Term{subgroup}{子群} $P_A$，\Term{smooth}{光滑}且\Term{fiber}{纤维}\Term{connected}{连通}，包含 $\operatorname{Centr}_G(Q)$，并满足
\[
 \operatorname{Lie}(P_A)=\operatorname{Lie}(G)^0\oplus\bigoplus_{\alpha\in R_A}\operatorname{Lie}(G)^\alpha.
\]
由 6.1、\Term{minimal-parabolic}{极小抛物子群}的\Term{conjugacy}{共轭}性，以及 $G$ 中包含 $Q$ 的\Term{parabolic-subgroup}{抛物子群}的\Term{root}{根}集是闭集这一事实（由 1.4 经\Term{split-group}{分裂}立即得到），立即推出：
\begin{SGAresult}{命题}{7.7}
（i）映射 $A\mapsto P_A$ 是从 $\Delta$ 的\Term{subset}{子集}的\Term{set}{集合}到 $G$ 中包含 $P$ 的\Term{parabolic-subgroup}{抛物子群}的\Term{set}{集合}的\Term{bijection}{双射}。此\Term{bijection}{双射}保持自然的包含\Term{order-relation}{序关系}。
（ii）$G$ 的每个\Term{parabolic-subgroup}{抛物子群}都由 $G(S)$ 的一个\Term{section}{截面}\Term{conjugacy}{共轭}到唯一的 $P_A$。
\end{SGAresult}

\textbf{7.8。}设 $P\supset Q$ 如上。考虑 $G$ 在 $S$ 上的相对\Term{root-datum}{根数据}（7.5）及\Term{canonical-isomorphism}{典范同构}
\[
 f:(M,M^*,R,R^*,R_+)\xrightarrow{\sim}
 (\underline M,\underline M^*,\underline R,\underline R^*,\underline R_+).
\]
$R_+$ 的\Term{simple-roots}{单根集} $\Delta$ 被变为 $\underline R_+$ 的\Term{simple-roots}{单根集} $\underline\Delta$，因此 $\Delta$ 的每个\Term{subset}{子集} $A$ 被变为\Term{subset}{子集} $f(A)\subset\underline\Delta$。
\begin{SGAresult}{定义}{7.9.0}
\footnote{编者注：为突出“\Term{relative-type}{相对型}”的定义，增添了编号 7.9.0。}设 $H$ 是 $G$ 的任意\Term{parabolic-subgroup}{抛物子群}。由 7.7（ii），它与唯一的 $P_A$ \Term{conjugacy}{共轭}。记 $t_r(H)=f(A)\subset\underline\Delta$。利用\Term{conjugacy}{共轭}定理立即可验，$t_r(H)$ 与对子 $P\supset Q$ 的选择无关。称它为 $H$ 的\Term{relative-type}{相对型}。
\end{SGAresult}
\SourcePage{327}
\begin{SGAresult}{命题}{7.9}
（i）映射 $H\mapsto t_r(H)$ 诱导 $G$ 的\Term{parabolic-subgroup}{抛物子群}在 $G(S)$ 下的\Term{conjugacy}{共轭}类的\Term{set}{集合}与 $\underline\Delta$ 的\Term{subset}{子集}的\Term{set}{集合}之间的\Term{bijection}{双射}。

（ii）设 $H$ 是 $G$ 的\Term{parabolic-subgroup}{抛物子群}，$P$ 是包含在 $H$ 中的一个\Term{minimal-parabolic}{极小抛物子群}，$Q$ 是 $P$ 的一个\Term{levi-subgroup}{Levi 子群}的最大\Term{center}{中心}\Term{split-torus}{分裂环面}，$\Delta$ 是 $P$ 相对于 $Q$ 的\Term{simple-roots}{单根集}，$f:\Delta\xrightarrow{\sim}\underline\Delta$ 是\Term{canonical-isomorphism}{典范同构}。则对每个 $\alpha\in\Delta$，有\Term{equivalence-relation}{等价关系}
\[
 f(\alpha)\in t_r(H)\quad\Longleftrightarrow\quad\operatorname{Lie}(H)^{-\alpha}\neq0,
\]
并且 $H=P_A$，其中 $A=f^{-1}(t_r(H))$。

（iii）若 $H$ 和 $H'$ 是 $G$ 中包含 $P$ 的两个\Term{parabolic-subgroup}{抛物子群}，则（另一些\Term{equivalence}{等价}条件见 3.8.1（ii）及 5.5（i））
\[
 t_r(H)\subset t_r(H')\quad\Longleftrightarrow\quad t(H)\subset t(H').
\]
\end{SGAresult}

\textbf{7.10。}可以研究两个\Term{minimal-parabolic}{极小抛物子群}的相对位置；结果如下（记号 $t_2(P,P_1)$ 见 4.5.2）：

（1）若 $P,P_1,P',P'_1$ 是 $G$ 的四个\Term{minimal-parabolic}{极小抛物子群}，则 $t_2(P,P_1)=t_2(P',P'_1)$（即 $(P,P_1)$ 和 $(P',P'_1)$ 对（fpqc）\Term{topology}{拓扑}\Term{local-scheme}{局部}\Term{conjugacy}{共轭}），当且仅当存在 $g\in G(S)$，使 $\operatorname{int}(g)P=P'$、$\operatorname{int}(g)P_1=P'_1$。

（2）特别地，固定一个以 $L$ 为\Term{levi-subgroup}{Levi 子群}的\Term{minimal-parabolic}{极小抛物子群} $P$，设 $T$ 是 $L$ 的一个\Term{maximal-torus}{极大环面}。考虑 $G$ 中相对于 $P$ 处于\Term{standard-position}{标准位置}的\Term{minimal-parabolic}{极小抛物子群}的\Term{scheme}{概形} $\operatorname{Par}_{t_{\min}}(G;P)$。有\Term{morphism}{态射}（见 4.5.5）
\[
 f:\operatorname{Par}_{t_{\min}}(G;P)\longrightarrow W_P(T)\backslash W_G(T)/W_P(T),
\]
其\Term{fiber}{纤维}是“$P$ 在 $\operatorname{Par}_{t_{\min}}(G;P)$ 中的\Term{orbit}{轨道}”。由（1），$f$ 因此诱导\Term{monomorphism}{单态射}
\[
 P(S)\backslash\operatorname{Par}_{t_{\min}}(G;P)(S)
 \hookrightarrow\bigl(W_P(T)\backslash W_G(T)/W_P(T)\bigr)(S).
\]
此\Term{morphism}{态射}的像可等同于 $W$；这就是 Bruhat 定理：$P(S)$ 在 $\operatorname{Par}_{t_{\min}}(G;P)(S)$ 中的每个\Term{orbit}{轨道}包含且只包含一个以 $L$ 为\Term{levi-subgroup}{Levi 子群}的 $G$ 的\Term{parabolic-subgroup}{抛物子群}（即形如 $\operatorname{int}(n)P$，其中 $n\in\operatorname{Norm}_G(L)$）。

（3）换言之，设 $E(S)$ 是使 $\operatorname{int}(g)P$ 和 $P$ 处于相互\Term{standard-position}{标准位置}的 $g\in G(S)$ 所成的\Term{set}{集合}。则 $E(S)$ 有一个按 $W$ 编号的模 $P(S)$ \Term{double-coset}{双陪集}分拆：
\[
 E(S)=P(S)\,W\,P(S)
\]
（记号意义显然）。置 $U=\operatorname{rad}^{u}(P)(S)$，也可写成
\[
 E(S)=U(S)\cdot\operatorname{Norm}_G(L)(S)\cdot U(S),
\]
从而显示出 $E(S)$ 的一个模 $U(S)$ 的\Term{double-coset}{双陪集}分拆，按 $\operatorname{Norm}_G(L)(S)$ 编号。

（4）若 $S$ 是一个域的谱，则 $E(S)=G(S)$，从而得到 [BT65]，5.15。

\SourcePage{328}
\textbf{反例 7.11。}设 $S=\operatorname{Spec}(\mathbb Z/4\mathbb Z)$，$G=\mathrm{SL}_{2,S}$。设 $B$ 是通常的\Term{borel-subgroup}{Borel 子群}，由满足 $c=0$ 的矩阵 $\left(\begin{smallmatrix}a&b\\c&d\end{smallmatrix}\right)$ 组成。设 $g=\left(\begin{smallmatrix}-1&1\\2&1\end{smallmatrix}\right)\in G(S)$，置 $B'=\operatorname{int}(g)B$。则 $B(S)=B'(S)$，而 $B\cap B'$ 不含\Term{maximal-torus}{极大环面}\footnote{编者注：事实上，$B'$ 是由方程 $c=2(a+b)$ 定义的 $G$ 的\Term{subgroup}{子群}；于是 $B\cap B'$ 在 $S$ 上不\Term{flat}{平坦}，故由 4.5.1 不含\Term{maximal-torus}{极大环面}。}。这一方面说明，两个不同的\Term{minimal-parabolic}{极小抛物子群}可以有相同的\Term{section}{截面}群；另一方面说明，一般不存在仅借助群 $P(S)$ 和 $P'(S)$ 判断两个\Term{minimal-parabolic}{极小抛物子群} $P$ 和 $P'$ 是否处于\Term{standard-position}{标准位置}的判据。特别地，$G(S)$ 的\Term{subset}{子集} $E(S)$ 似乎不能仅借助配置 $\{G(S),P(S),\operatorname{Norm}_G(L)(S)\}$ 来定义（在前例中，这个\Term{subset}{子集}由 $c\neq2$ 定义\footnote{编者注：更一般地，对每个 $S$ \Term{scheme}{概形} $S'$，$E(S')$ 由条件“$c$ 为零或\Term{invertible}{可逆}”定义。}）。

\textbf{7.12。}现在拟研究 $\mathscr R(G/S)$ 随 $S$ 的变化。设 $S'$ 是一个也半\Term{local-scheme}{局部}、\Term{connected}{连通}、\Term{nonempty}{非空}的 $S$ \Term{scheme}{概形}。取 $G$ 的一个极大\Term{split-torus}{分裂环面} $Q$；则 $Q_{S'}$ 是 $G_{S'}$ 的\Term{split-torus}{分裂环面}。设 $Q'$ 是 $G_{S'}$ 中包含 $Q_{S'}$ 的极大\Term{split-torus}{分裂环面}。置
\begin{align*}
 M&=\operatorname{Hom}_{S\text{-gr}}(Q,\mathbb G_{m,S})\simeq\operatorname{Hom}_{S'\text{-gr}}(Q_{S'},\mathbb G_{m,S'}),\\
 M'&=\operatorname{Hom}_{S'\text{-gr}}(Q',\mathbb G_{m,S'}).
\end{align*}
\Term{monomorphism}{单态射} $Q_{S'}\to Q'$ 诱导满\Term{homomorphism}{同态} $u:M'\to M$。记
\[
 L=\operatorname{Centr}_G(Q),\qquad L'=\operatorname{Centr}_{G_{S'}}(Q'),
\]
则 $L'\subset L_{S'}$。
若 $H$ 是 $G$ 中包含 $L$ 的\Term{subgroup}{子群}，则 $H_{S'}$ 包含 $L'$，且有
\begin{align*}
 \operatorname{Lie}(H)&=\operatorname{Lie}(L)\oplus\bigoplus_{\alpha\in R_H}\operatorname{Lie}(H)^\alpha,\\
 \operatorname{Lie}(H_{S'})&=\operatorname{Lie}(L')\oplus\bigoplus_{\alpha'\in R'_{H_{S'}}}\operatorname{Lie}(H_{S'})^{\alpha'},
\end{align*}
其中 $R_H$（分别为 $R'_{H_{S'}}$）表示 $H$（分别为 $H_{S'}$）相对于 $Q$（分别为 $Q'$）的\Term{root}{根}集。立即得到
\[
 R_H\subset u(R'_{H_{S'}})\subset R_H\cup\{0\}.
\]
取 $H=G$，首先看到 $R\subset u(R')\subset R\cup\{0\}$；再取 $H$ 为以 $L$ 为\Term{levi-subgroup}{Levi 子群}的\Term{minimal-parabolic}{极小抛物子群} $P$，则 $R'_{H_{S'}}$ 包含 $R'$ 的一个\Term{positive-roots}{正根系}，故由 7.4，存在 $G_{S'}$ 中以 $L'$ 为\Term{levi-subgroup}{Levi 子群}且包含在 $P_{S'}$ 中的\Term{minimal-parabolic}{极小抛物子群} $P'$。
\SourcePage{329}
因此构造出图
\[
\begin{tikzcd}
 Q_{S'}\arrow[r,hook]\arrow[d,hook]&L_{S'}\arrow[r,hook]&P_{S'}\\
 Q'\arrow[r,hook]&L'\arrow[u,hook]\arrow[r,hook]&P'\arrow[u,hook].
\end{tikzcd}
\]
若 $R_+$（分别为 $\Delta$）是由 $P$ 定义的 $R$ 的\Term{positive-roots}{正根系}（分别为\Term{simple-roots}{单根系}），同样定义 $R'_+$ 和 $\Delta'$，则容易验证
\[
 R_+\subset u(R'_+)\subset R_+\cup\{0\},\qquad
 \Delta\subset u(\Delta')\subset\Delta\cup\{0\}.
\]
现在设 $w\in W\simeq\operatorname{Norm}_G(Q)(S)/\operatorname{Centr}_G(Q)(S)$，取 $w$ 的代表元 $n\in\operatorname{Norm}_G(Q)(S)$。有 $\operatorname{int}(n)Q=Q$，故 $\operatorname{int}(n)L=L$，从而 $\operatorname{int}(n)L_{S'}=L_{S'}$。于是 $Q'$ 和 $\operatorname{int}(n)Q'$ 是 $L_{S'}$ 的两个极大\Term{split-torus}{分裂环面}，故由\Term{section}{截面} $x\in L(S')$ \Term{conjugacy}{共轭}，且有 $\operatorname{int}(nx)Q'=Q'$，从而 $nx\in\operatorname{Norm}_{G_{S'}}(Q')(S')$。设 $w'$ 是 $n'=nx$ 在 $W'\simeq\operatorname{Norm}_{G_{S'}}(Q')(S')/\operatorname{Centr}_{G_{S'}}(Q')(S')$ 中的像。显然，$w'$ 在 $M'$ 上的\Term{action}{作用}与\Term{projection}{投影} $u:M'\to M$ 相容，并且在 $M$ 上诱导的\Term{action}{作用}与 $w$ 定义的\Term{action}{作用}相同。
现在利用相对\Term{root-datum}{根数据}的定义及\Term{conjugacy}{共轭}定理，不难证明：
\begin{SGAresult}{定理}{7.13}
设 $S$ 和 $S'$ 是两个\Term{nonempty}{非空}\Term{connected}{连通}\Term{semilocal-scheme}{半局部概形}，$S'\to S$ 是 $S$ \Term{scheme}{概形}的\Term{morphism}{态射}，$G$ 是\Term{reductive-group}{约化 $S$ 群}，
\begin{align*}
 \mathscr R(G/S)&=(\underline M,\underline M^*,\underline R,\underline R^*,\underline R_+),\\
 \mathscr R(G_{S'}/S')&=(\underline M',\underline M'^*,\underline R',\underline R'^*,\underline R'_+)
\end{align*}
是相对\Term{pinned-root-datum}{带钉扎根数据}。存在\Term{canonical-homomorphism}{典范同态}
\[
 \underline u:\underline M'\longrightarrow\underline M,
\]
满足下列条件：
（i）$\underline u$ \Term{surjectivity}{满射}。
（ii）对每个 $w\in W$，存在 $W'$ 的元素 $w'$，与 $\underline u$ 相容，且在 $\underline M$ 上诱导 $w$。
（iii）对 $\underline M$ 的每个\Term{subset}{子集} $X$，记 $X^\wedge=X\cap(\underline M-\{0\})$。则
\[
 \underline u(\underline R'_+)^\wedge=\underline R_+,\qquad
 \underline u(\underline\Delta')^\wedge=\underline\Delta.
\]
（iv）对每个\Term{parabolic-subgroup}{抛物子群} $H$，考虑 $t_r(H)\subset\underline\Delta$ 和 $t_r(H_{S'})\subset\underline\Delta'$。则
\begin{align*}
 t_r(H_{S'})
 &=\underline u^{-1}\bigl(t_r(H)\cup\{0\}\bigr)\cap\underline\Delta'\\
 &=\{\alpha'\in\underline\Delta'\mid\underline u(\alpha')\in t_r(H)\text{ 或 }\underline u(\alpha')=0\}.
\end{align*}
\end{SGAresult}
\textbf{注记 7.14。}若 $G$ 在 $S$ 上可\Term{split-group}{分裂}，它的极大\Term{split-torus}{分裂环面}就是\Term{maximal-torus}{极大环面}，因而这里引进的相对概念与此前引进的绝对记号一致。因此，前一定理借助群 $G_{S'}$ 的绝对\Term{root-datum}{根数据}与绝对型，给出相对\Term{root-datum}{根数据} $\mathscr R(G/S)$ 和\Term{relative-type}{相对型} $t_r$ 的描述，其中选择 $S'$ 使 $G_{S'}$ 可\Term{split-group}{分裂}（见第 XXIV 讲 4.4.1）。关于这种描述，参见 [BT65]，6.12 及以下各段。

\SourcePage{330}
\textbf{7.15。}设 $S$ 是 Hensel \Term{local-scheme}{局部}\Term{scheme}{概形}，$s_0$ 是其\Term{closed-point}{闭点}，$S_0$ 是 $s_0$ 的\Term{residue-field}{剩余域}的谱，等同于 $S$ 的一个\Term{closed-subscheme}{闭子概形}；对 $S$ 上每个\Term{object}{对象} $X$，以 $X_0$ 表示由 $X$ 作\Term{base-change}{基变换}得到的 $S_0$ 上的\Term{object}{对象}。最后设 $G$ 是\Term{reductive-group}{约化 $S$ 群}。对 $G$ 的每个\Term{parabolic-subgroup}{抛物子群} $P$，$P_0$ 是 $G_0$ 的\Term{parabolic-subgroup}{抛物子群}；反之，对 $G_0$ 的每个\Term{parabolic-subgroup}{抛物子群} $\overline P$，存在 $G$ 的\Term{parabolic-subgroup}{抛物子群} $P$，使 $P_0=\overline P$（这由 Hensel 引理及 $\operatorname{Par}(G)$ 是\Term{smooth}{光滑 $S$ 概形}得到）。特别地（见 5.7），$G$ 的\Term{parabolic-subgroup}{抛物子群} $P$ 是\Term{minimal-parabolic}{极小的}，当且仅当 $P_0$ 是\Term{minimal-parabolic}{极小的}。选定 $G$ 的这样一个\Term{subgroup}{子群} $P$ 后，类似的论证说明，$P_0$ 的极大\Term{split-torus}{分裂子环面}都形如 $T_0$，其中 $T$ 是 $P$ 的极大\Term{split-torus}{分裂子环面}。不难推出，$G$ 在 $S$ 上和 $G_0$ 在 $S_0$ 上的相对\Term{root-datum}{根数据}\Term{canonical-isomorphism}{典范同构}，所以 $G$ 的\Term{parabolic-subgroup}{抛物子群}理论归结为 $G_0$ 的\Term{parabolic-subgroup}{抛物子群}理论。
还须指出，每个\Term{reductive-group}{约化 $S_0$ 群}都形如 $G_0$（第 XXIV 讲，命题 1.21），这又反过来将一个\Term{reductive-group}{约化 $S_0$ 群}的\Term{parabolic-subgroup}{抛物子群}的研究归结为 $S$ 上的相应研究。

\textbf{参考文献}
\begin{description}
\item[{[BT65]}] A. Borel，J. Tits，\textit{\Term{reductive-group}{约化群}}（Groupes réductifs），Publ. Math. I.H.É.S. \textbf{27}（1965），55--150。\footnote{编者注：原版仅有此条参考文献，我们增添了以下各条。}
\item[{[BT71]}] A. Borel，J. Tits，\textit{\Term{reductive-group}{约化群}的\Term{unipotent}{幺幂}元素与\Term{parabolic-subgroup}{抛物子群}，I}（Éléments unipotents et sous-groupes paraboliques de groupes réductifs. I），Invent. Math. \textbf{12}（1971），95--104。
\item[{[Ch05]}] C. Chevalley，\textit{\Term{semisimple}{半单}\Term{algebraic-group}{代数群}的分类}（Classification des groupes algébriques semi-simples；与 P. Cartier、A. Grothendieck、M. Lazard 合作），Collected Works，第 3 卷，Springer，2005。
\item[{[DG70]}] M. Demazure，P. Gabriel，\textit{\Term{algebraic-group}{代数群}}（Groupes algébriques），Masson \& North-Holland，1970。
\item[{[Gi71]}] J. Giraud，\textit{\Term{nonabelian-cohomology}{非阿贝尔上同调}}（Cohomologie non abélienne），Springer-Verlag，1971。
\end{description}
